\documentclass[11pt,a4paper]{article}

% ---------------------------------------------------------------------------
% Packages
% ---------------------------------------------------------------------------
\usepackage[utf8]{inputenc}
\usepackage[T1]{fontenc}
\usepackage{amsmath}
\usepackage[margin=1in]{geometry}
\usepackage{booktabs}
\usepackage{enumitem}
\usepackage{xcolor}
\usepackage{listings}
\usepackage{longtable}
\usepackage{titlesec}
\usepackage[hidelinks]{hyperref}

% ---------------------------------------------------------------------------
% Listings (code) style
% ---------------------------------------------------------------------------
\definecolor{codebg}{gray}{0.96}
\lstset{
  basicstyle=\ttfamily\small,
  breaklines=true,
  columns=fullflexible,
  frame=single,
  framesep=4pt,
  backgroundcolor=\color{codebg},
  showstringspaces=false,
  keepspaces=true,
}

\titleformat{\section}{\large\bfseries}{\thesection.}{0.6em}{}
\setlist{itemsep=2pt,topsep=3pt}

% ---------------------------------------------------------------------------
% Title
% ---------------------------------------------------------------------------
\title{\textbf{Human Safety Monitoring System}\\[2pt]
\large PPE Detection and Person Tracking --- Technical Report}
\author{Arif Faishal}
\date{23 June 2026}

\begin{document}
\maketitle

\begin{abstract}
\noindent
This report documents a computer-vision service for workplace human-safety
monitoring. The system combines a fine-tuned YOLO model for Personal
Protective Equipment (PPE) detection with a pretrained YOLOv11 model for
person detection and multi-object tracking. It exposes a FastAPI HTTP
interface offering single-frame detection (PPE or person) and full-video
person tracking, where each tracked person is annotated with their PPE
compliance state (helmet, vest, mask). Tracking-identity consistency is
quantified through an ID Switch Rate metric, and throughput is reported as
frames per second (FPS). The document covers the data, models, evaluation
methodology, system architecture, the primary AI-assisted development steps,
and current limitations.
\end{abstract}

% ===========================================================================
\section{Data}
% ===========================================================================

\subsection{Inputs}
The service operates on two media modalities supplied by the client at request
time:
\begin{itemize}
  \item \textbf{Images} (JPEG / PNG) --- used by the single-frame detection
        endpoint. Bytes are decoded to a BGR array with OpenCV
        (\verb|cv2.imdecode|).
  \item \textbf{Video} (MP4) --- used by the person-tracking endpoint. The
        upload is buffered to a temporary file and read frame-by-frame with
        \verb|cv2.VideoCapture|.
\end{itemize}
A representative operational sample, \verb|datas/hardhat.mp4|, is retained for
local testing alongside generated artifacts under \verb|datas/detections/|.

\subsection{PPE class taxonomy}
The PPE model emits six mutually paired classes (a positive ``wearing'' class
and its ``missing'' counterpart) that are mapped onto the application's
\verb|SafetyLabel| enumeration. The class index ordering matches the project
reference in \verb|scripts/tracker.py|.

\begin{center}
\begin{tabular}{cll}
\toprule
\textbf{Index} & \textbf{Model class} & \textbf{SafetyLabel} \\
\midrule
0 & Hard Hat       & \verb|hard_hat| \\
1 & No Hard Hat    & \verb|no_hard_hat| \\
2 & Mask           & \verb|mask| \\
3 & No Mask        & \verb|no_mask| \\
4 & Safety Vest    & \verb|safety_vest| \\
5 & No Safety Vest & \verb|no_safety_vest| \\
\bottomrule
\end{tabular}
\end{center}

\noindent
The person model is class-agnostic with respect to PPE; only COCO class
\verb|0| (\emph{person}) is requested at inference. Detected persons are mapped
to the \verb|person| label.

\subsection{Outputs}
Two artifact families are persisted to the configured output directory
(\verb|DETECTION_OUTPUT_DIR|, default \verb|datas/detections|):
\begin{itemize}
  \item \textbf{Detection}: an annotated image
        (\verb|{type}_detection_{id}.{ext}|) plus a JSON response
        (\verb|.json|) carrying every bounding box, label, and confidence.
  \item \textbf{Tracking}: an annotated MP4 (\verb|track_person_{id}.mp4|) and a
        per-frame JSON file (\verb|track_person_{id}.json|) containing tracked
        persons, their PPE booleans, the ID Switch Rate, and the average FPS.
\end{itemize}

% ===========================================================================
\section{Model}
% ===========================================================================

The system wraps two YOLO models through the Ultralytics runtime
(\verb|ultralytics>=8.4.75|). Both are loaded once during application startup
and shared as a single in-memory instance for all requests.

\begin{center}
\begin{tabular}{lll}
\toprule
\textbf{Role} & \textbf{Weights} & \textbf{Size} \\
\midrule
PPE detection            & \verb|models/yolo_50.pt|             & $\approx$161\,MB \\
Person detection/tracker & \verb|models/yolo_person_tracker.pt| & $\approx$5.6\,MB \\
\bottomrule
\end{tabular}
\end{center}

\begin{itemize}
  \item \textbf{PPE model} (\verb|yolo_50.pt|): a YOLO detector fine-tuned on
        the six-class PPE taxonomy above. Its larger footprint is consistent
        with a medium/large backbone trained for the construction-safety
        domain.
  \item \textbf{Person model} (\verb|yolo_person_tracker.pt|): a pretrained
        YOLOv11 model (nano-scale, given its size) used both for person
        detection and as the tracking detector.
\end{itemize}

\subsection{Loading lifecycle}
Model weights are loaded inside the FastAPI \verb|lifespan| context manager so
that the cost is paid once at boot, never on the request path. A readiness
guard (\verb|is_ready()|) requires \emph{both} models to be present before any
endpoint will serve traffic; otherwise the API returns
\verb|503 Service Unavailable|. The \verb|/ready| probe surfaces this state for
orchestration.

\subsection{Inference configuration}
\begin{itemize}
  \item \textbf{Confidence threshold}: a single configurable value
        (\verb|PERSON_DETECTION_THRESHOLD|, default \verb|0.3|) gates both
        detection and tracking, avoiding hard-coded thresholds.
  \item \textbf{Tracking}: ByteTrack (\verb|bytetrack.yaml|) is used as the
        association algorithm. Tracking is driven frame-by-frame with
        \verb|persist=True| so each model maintains its own track state across
        the video.
\end{itemize}

% ===========================================================================
\section{Evaluation}
% ===========================================================================

Evaluation focuses on the two operational qualities that matter for a live
monitoring service: \emph{tracking-identity stability} and \emph{throughput}.

\subsection{Tracking-ID consistency --- ID Switch Rate}
Identity stability is measured with an ID Switch Rate. A spatial cell is
defined by quantising a bounding box to a 50\,px grid; whenever the track ID
occupying a previously seen cell changes between observations, an identity
switch is counted. The rate is
\[
\text{ID Switch Rate} \;=\; \frac{\#\,\text{ID switches}}{\#\,\text{frames}} \times 100\%.
\]
This is reported per tracking run in the JSON output, e.g.
\verb|{"id_switches": 0, "total_frames": 25, "rate": "0.00%"}|. The method is a
lightweight, ground-truth-free proxy for identity continuity (in the spirit of
MOT identity metrics) rather than a formal IDF1/MOTA computation.

\subsection{Throughput --- FPS}
Per-frame processing latency is measured and converted to an instantaneous FPS
value, which is overlaid on the top-left of every annotated frame. The average
FPS over the clip is reported in the summary. In a CPU smoke test on
$640\times360$ clips, the end-to-end pipeline (two YOLO models plus ByteTrack
per frame) ran at roughly \textbf{9\,FPS}.

\subsection{Functional verification}
Each module was verified end-to-end via FastAPI's \verb|TestClient|:
\begin{itemize}
  \item Both models load on boot and \verb|/ready| reports ready.
  \item \verb|/detect| returns \verb|200| for \verb|ppe| and \verb|person|, and
        \verb|422| for an invalid \verb|detection_type|.
  \item \verb|/track-person| returns \verb|200|; track IDs, per-person PPE
        booleans, the ID Switch Rate, FPS, and both output files are produced
        (17 person-detections across a 25-frame test segment).
\end{itemize}

\subsection{Not yet measured}
No labelled ground truth is wired in, so detection quality
(precision/recall/mAP) and formal tracking accuracy (MOTA/IDF1) are not
reported. These require an annotated evaluation set and are noted in
Limitations.

% ===========================================================================
\section{System Design}
% ===========================================================================

The application follows a layered FastAPI architecture with strict separation
between HTTP concerns and model logic.

\begin{lstlisting}
app/
  main.py                 # App factory + lifespan (loads both models on boot)
  core/
    config.py             # pydantic-settings (paths, thresholds, output dir)
    logging.py            # Logging setup
  api/
    deps.py               # Shared singletons (get_detector, get_monitor)
    v1/
      router.py           # Aggregates all v1 routers
      endpoints/
        health.py         # /health, /ready
        detections.py     # /detections (legacy single-frame)
        monitoring.py     # /monitoring/evaluate, /monitoring/alerts
        detect.py         # /detect  (detection_type: ppe | person)
        track.py          # /track-person (video -> annotated video + JSON)
  schemas/                # Pydantic DTOs (detection, monitoring, tracking)
  services/
    detector.py           # YOLO wrapper: load, detect, annotate, track_persons
    safety_monitor.py     # Detections -> events -> alerts (framework-agnostic)
  utils/
    storage.py            # Persistence helpers (image/JSON/video artifacts)
\end{lstlisting}

\subsection{Layers}
\begin{itemize}
  \item \textbf{API layer} (\verb|api/v1/endpoints|): request validation,
        status codes, and response models. No model code lives here.
  \item \textbf{Service layer} (\verb|services/detector.py|): owns both models
        privately and exposes a stable interface --- \verb|detect()|,
        \verb|annotate()|, and \verb|track_persons()|. Free of FastAPI imports.
  \item \textbf{Schema layer} (\verb|schemas/|): typed DTOs returned by every
        endpoint (e.g.\ \verb|PersonDetectionResult|, \verb|PersonTrackingResult|,
        \verb|TrackingData|).
  \item \textbf{Utility layer} (\verb|utils/storage.py|): file-naming templates
        and artifact persistence, decoupled from the request handlers.
  \item \textbf{Core} (\verb|core/|): centralised configuration and logging.
\end{itemize}

\subsection{Person-tracking pipeline}
For each frame of the uploaded video:
\begin{enumerate}
  \item Run the person model with ByteTrack (\verb|classes=[0]|, \verb|persist=True|)
        to obtain tracked person boxes and IDs.
  \item Run the PPE model with ByteTrack to obtain PPE class boxes.
  \item Associate PPE to a person when a PPE box \emph{centre} falls inside the
        person box, deriving the booleans \verb|helmet|, \verb|vest|, \verb|mask|.
  \item Draw each person box with the meta label
        \verb|ID{n} H:{0/1} V:{0/1} M:{0/1}| (green if fully compliant, red
        otherwise) and overlay the current FPS.
  \item Update the ID Switch Rate accumulator and append the frame record.
\end{enumerate}
The annotated frames are written to an MP4 (\verb|mp4v|), and the per-frame
records plus summary metrics are serialised to JSON.

\subsection{Request example}
\begin{lstlisting}[language=bash]
curl -X POST http://localhost:8000/api/v1/track-person \
  -F "file=@datas/hardhat.mp4"
\end{lstlisting}

% ===========================================================================
\section{AI Usage Log (Primary)}
% ===========================================================================

The system was built incrementally with AI assistance; every session is
recorded in \verb|changelogs/agent-changelogs.md|. The five \emph{primary}
feature-delivering steps are summarised below.

\begin{longtable}{p{0.5cm} p{4.3cm} p{9.0cm}}
\toprule
\textbf{\#} & \textbf{Prompt (paraphrased)} & \textbf{Outcome} \\
\midrule
\endhead
3 & Add a boot-time model loader for the PPE and person weights. &
Added \verb|PPE_MODEL_PATH| / \verb|PERSON_MODEL_PATH| settings; loads both YOLO
models in the lifespan; \verb|is_ready()| now requires both. \\
\addlinespace
4 & Add a \verb|detect-person| endpoint (image upload, threshold 0.3). &
Implemented person detection (COCO class 0) returning
\verb|PersonDetectionResult|; added \verb|PERSON_DETECTION_THRESHOLD| and
\verb|DETECTION_OUTPUT_DIR|. \\
\addlinespace
5 & Add the PPE detection module and align labels to the class taxonomy. &
Replaced placeholder labels with the six PPE classes; added the class-index
mapping and the \verb|detect-ppe| path. \\
\addlinespace
6 & Unify the two endpoints and save annotated artifacts. &
Merged into a single \verb|/detect| endpoint keyed by \verb|detection_type|;
added OpenCV annotation and paired image\,+\,JSON persistence. \\
\addlinespace
8 & Build the person-tracking module. &
Added ByteTrack person\,+\,PPE tracking, the ID Switch Rate, FPS overlay, and
the \verb|/track-person| endpoint producing an annotated video and tracking
JSON. \\
\bottomrule
\end{longtable}

% ===========================================================================
\section{Limitations}
% ===========================================================================

\begin{itemize}
  \item \textbf{Synchronous video processing.} Tracking runs inline within the
        HTTP request, holding the connection open for the duration of the
        video. Long clips should be moved to a background job with a
        poll/callback contract.
  \item \textbf{CPU throughput.} At $\approx$9\,FPS on CPU for $640\times360$
        input, the pipeline is not real-time; GPU execution or a smaller PPE
        backbone would be needed for live streams.
  \item \textbf{Approximate PPE--person association.} Attribution uses a simple
        ``centre-in-box'' test, which can mis-assign PPE when people overlap
        heavily. IoU-based or nearest-person matching would be more robust.
  \item \textbf{Heuristic identity metric.} The ID Switch Rate uses a 50\,px
        spatial-cell proxy, not a ground-truth alignment; it is indicative
        rather than a formal MOTA/IDF1 score.
  \item \textbf{No labelled evaluation.} Detection accuracy (precision, recall,
        mAP) and tracking accuracy are not yet measured against annotated data.
  \item \textbf{Conservative compliance semantics.} A PPE item is reported
        \verb|true| only when its positive class is detected near the person;
        missed detections are therefore treated as non-compliant, which can
        inflate false ``violation'' signals.
  \item \textbf{Independent trackers.} The person and PPE models maintain
        separate ByteTrack states, so PPE identities are not themselves linked
        to a person identity over time --- only re-associated per frame.
  \item \textbf{Output encoding.} Annotated video is written with the
        \verb|mp4v| codec and no audio; container/codec compatibility is not
        configurable.
\end{itemize}

\end{document}
